\documentclass[12pt,a4paper]{article}

\usepackage[utf8]{utf8}
\usepackage[margin=1in]{geometry}
\usepackage{amsmath,amssymb,amsfonts}
\usepackage{graphicx}
\usepackage{booktabs}
\usepackage{hyperref}
\usepackage{cite}
\usepackage{xcolor}
\usepackage{tikz}
\usetikzlibrary{shapes,arrows,positioning}

\definecolor{primaryColor}{RGB}{0, 119, 182}
\definecolor{secondaryColor}{RGB}{43, 147, 72}

\hypersetup{
    colorlinks=true,
    linkcolor=primaryColor,
    citecolor=secondaryColor,
    urlcolor=primaryColor
}

\title{\textbf{AquaFusion AI: Project Proposal}\\
\large Autonomous Aquaponics Management with Non-Feeder Plant-Based Biomass Cycling, Adaptive Overhead Spectrum Lighting, and Computer Vision Ecosystem Health Monitoring}
\author{\textbf{Samiru Nuwanaka} \\ Team SynchroCrop \\ Biomedical Engineering $\cdot$ Artificial Intelligence $\cdot$ Smart Aquaculture}
\date{\today}

\begin{document}

\maketitle

\begin{abstract}
Aquaponics presents an eco-friendly food production system, yet conventional implementations suffer from high operational labor, dependency on commercial fish feeds, energy-inefficient artificial lighting, and catastrophic ecosystem failures. This proposal presents \textbf{AquaFusion AI}, a novel, fully autonomous indoor aquaponics solution designed specifically for urban apartment/flat dwellers. The proposed system features: (1) \textbf{Computer Vision Fish Behavior Monitoring} using optical flow, trajectory tracking, and surface piping detection; (2) \textbf{Multimodal Health Diagnostics} for both aquatic life and plants (NDVI analysis, leaf chlorosis detection, ORP/DO sensors); (3) \textbf{Feederless Bio-Integrated Fish Nutrition} utilizing overhead cultivated aquatic plant beds (Duckweed/Lemna minor and Wolffia) for continuous direct grazing; and (4) \textbf{Adaptive Overhead Green-Shifted Lighting} that extracts maximum electrical power conversion efficiency, dissipates only photosynthetic active radiation (PAR) frequencies, and dynamic height/intensity adaptation to achieve \textbf{$>2\times$ crop yield} compared to standard windowsill/sunlight-restricted flat environments.
\end{abstract}

\tableofcontents
\newpage

\section{Introduction \& Problem Statement}
Urban residential flat dwellers face significant space, light, and logistical constraints when attempting sustainable food production:
\begin{enumerate}
    \item \textbf{Sunlight Deficit in Apartment Flats:} Windowsills and balconies provide intermittent, low-intensity natural sunlight, yielding sub-optimal photosynthetic rates and low crop growth.
    \item \textbf{Energy Inefficiency of Full-Spectrum Grow Lights:} Broad-spectrum lights waste power across unabsorbed wavelengths, generating excessive heat and high electricity bills in confined flat living spaces.
    \item \textbf{Commercial Feeder Dependence \& Maintenance Burden:} Automated mechanical feeders clog, require frequent replenishment, and introduce excessive protein waste leading to toxic ammonia spikes.
    \item \textbf{Ecosystem Fragility:} Sudden changes in dissolved oxygen ($\text{DO}$), pH, or ammonia ($\text{NH}_3$) can cause silent fish mass mortality before manual detection.
\end{enumerate}

\textbf{AquaFusion AI} solves these challenges by combining AI edge computer vision, closed-loop aquatic plant feeding, and adaptive PAR lighting designed specifically for indoor urban agriculture.

\section{Computer Vision Fish Behavior Monitoring}
To eliminate manual inspection, an underwater RGB-D camera connected to an edge microcontroller (Raspberry Pi 5 / Jetson Orin Nano) continuously analyzes fish swimming kinetics.

\subsection{Behavioral Metrics and Trajectories}
The system tracks three principal behavioral indicators:
\begin{itemize}
    \item \textbf{Swarming \& Spatial Distribution:} Normal fish disperse homogeneously across the water volume. Aggregation near the surface indicates hypoxia ($\text{DO} < 4.0\text{ mg/L}$), while bottom crowding indicates thermal stress or toxic ammonia spikes ($\text{NH}_3 > 0.05\text{ mg/L}$).
    \item \textbf{Velocity Vector Profiling (Optical Flow):} Using Lucas-Kanade optical flow, normal cruising speed $\vec{v}_f$ is calculated. Lethargy ($\|\vec{v}_f\| \rightarrow 0$) or hyperactive darting ($\|\vec{v}_f\| \gg v_{\text{baseline}}$) triggers immediate alert protocols.
    \item \textbf{Surfacing / Piping Detection:} Frequent vertical movement toward the surface air-water interface is detected using YOLOV8-nano object tracking, signaling emergency oxygen depletion.
\end{itemize}

\section{Investigating Fish \& Plant Health}

\subsection{Fish Health Diagnostics}
Fish health is evaluated dynamically via non-invasive multimodal sensing:
\begin{itemize}
    \item \textbf{Symptom Detection:} Image segmentation models detect body lesions, white spots (\textit{Ichthyophthirius multifiliis}), fin rot, and abnormal spinal curvature.
    \item \textbf{Water Quality Correlation:} Real-time chemical telemetry—Oxidation-Reduction Potential (ORP), Dissolved Oxygen (DO), pH, Electrical Conductivity (EC), Ammonia ($\text{NH}_3$), Nitrite ($\text{NO}_2^-$), and Nitrate ($\text{NO}_3^-$)—provides predictive root-cause diagnosis.
\end{itemize}

\subsection{Plant Health Diagnostics}
Plant foliage health is evaluated using top-down camera arrays:
\begin{itemize}
    \item \textbf{Foliar Chlorosis \& Deficiencies:} Computer vision monitors leaf color gradients. Yellowing older leaves indicate Nitrogen ($\text{N}$) deficiency, interveinal chlorosis indicates Iron ($\text{Fe}$) or Magnesium ($\text{Mg}$) deficiency, and purple tinting indicates Phosphorus ($\text{P}$) deficiency.
    \item \textbf{Growth Rate Estimation:} 3D depth sensors measure canopy volume expansion over 24-hour cycles.
\end{itemize}

\section{Feederless Aquatic Plant-Based Feeding System}
Traditional fish feeders are replaced by an integrated, self-sustaining aquatic plant foraging bed:
\begin{itemize}
    \item \textbf{Live Bio-Feeder Bed:} High-protein fast-growing aquatic macrophytes (\textit{Lemna minor} / Duckweed and \textit{Wolffia globosa}) are cultivated in a shallow sub-chamber directly above or accessible to the fish tank.
    \item \textbf{Direct Grazing \& Bio-Ramp:} Omnivorous and herbivorous fish (e.g., Tilapia, Koi, Gourami) continuously feed on fresh biomass enriched with plant-synthesized vitamins and protein ($35-45\%$ crude protein content).
    \item \textbf{Ammonia Recycling Loop:} Fish waste directly fertilizes the forage crop bed, removing toxic ammonia while producing fresh forage without human intervention or mechanical feeder jams.
\end{itemize}

\section{Adaptive Overhead Green-Shifted PAR Lighting Mechanism}

\subsection{Resolving Flat-Dweller Sunlight Limitations (More Than 2x Yield)}
In dense apartment flats, natural light intensity is typically limited to $50-150\ \mu\text{mol}/\text{m}^2/\text{s}$ for less than 4 hours daily. \textbf{AquaFusion AI} delivers optimal Photosynthetically Active Radiation (PAR) continuously.

\subsection{Narrow-Spectrum PAR Frequency Optimization \& Heat Dissipation}
\begin{itemize}
    \item \textbf{Selective Wavelength Emission:} Plants utilize specific wavelengths ($450\text{nm}$ Blue, $525-550\text{nm}$ Green penetration, $660\text{nm}$ Deep Red). Rather than burning energy across unabsorbed spectra (which causes excessive thermal radiation), AquaFusion AI uses high-efficiency solid-state LEDs emitting \emph{only} exact photosynthetic frequencies.
    \item \textbf{Green Light Canopy Penetration:} Deep canopy leaves under shade are penetrated efficiently by tuned green spectrum wavelengths ($525-550\text{nm}$), enabling lower-tier leaves to perform photosynthesis at full capacity.
    \item \textbf{Power Extraction Efficiency:} By filtering out non-photosynthetic thermal infrared and UV waste, power conversion efficiency increases by $>60\%$, dissipating minimal heat into flat living quarters.
\end{itemize}

\subsection{Adaptive Height \& Growth Motorized Tracking}
\begin{itemize}
    \item To maintain optimum Photosynthetic Photon Flux Density (PPFD) without burning growing tips, the overhead lighting array is mounted on motorized linear actuators.
    \item As the computer vision system detects canopy growth ($\Delta h$), the light bed automatically adjusts its height $h_{\text{light}}(t) = h_{\text{canopy}}(t) + d_{\text{optimal}}$, optimizing photon capture efficiency throughout the plant life cycle.
\end{itemize}

\section{Expected Results \& Metrics}

\begin{table}[h!]
\centering
\caption{System Performance Comparisons}
\begin{tabular}{lcc}
\toprule
\textbf{Metric} & \textbf{Standard Flat Sunlight} & \textbf{AquaFusion AI System} \\
\midrule
Daily PAR Yield ($\text{mol}/\text{m}^2/\text{day}$) & $4.2 - 6.5$ & \textbf{$15.0 - 22.0$ ($>2\times$ Growth)} \\
Feeder Maintenance & Refill every 3-5 days & \textbf{Self-sustaining (0 Days)} \\
Power Conversion Loss & High (Broad spectrum heat) & \textbf{Minimal (Tuned PAR only)} \\
Fish Survival Rate & $75-80\%$ & \textbf{$>98\%$} \\
\bottomrule
\end{tabular}
\end{table}

\section{Conclusion}
\textbf{AquaFusion AI} redefines urban indoor aquaponics. By integrating computer vision fish behavior tracking, automated plant/fish disease diagnostics, a feederless aquatic plant grazing mechanism, and adaptive green-penetration PAR lighting, flat dwellers can achieve more than double the crop yield of traditional natural-light farming while eliminating daily maintenance and minimizing power consumption.

\end{document}
